\chapter{Was ist Heldenpfade?}

\section*{Hinweis zur Sprache}

Zur besseren Lesbarkeit wird in diesem Regelwerk überwiegend die Bezeichnung „Spieler“ verwendet. Gemeint sind dabei selbstverständlich immer alle Spielerinnen und Spieler. Gleiches gilt für Begriffe wie „Held“ oder „Spielleitung“, die unabhängig vom Geschlecht verwendet werden.

\section*{Die wichtigste Regel}

Heldenpfade soll allen am Spieltisch Spaß machen.

Die Regeln dieses Regelwerks bilden einen gemeinsamen Rahmen. Sie sollen dabei helfen, spannende Abenteuer zu erleben und gemeinsam zu entscheiden, wie die Geschichte weitergeht.

Wenn ihr beim Spielen merkt, dass eine andere Lösung oder Regel eurer Gruppe mehr Spaß macht, dann ändert diese Regel einfach. Die Spielleitung entscheidet gemeinsam mit den Spielern, welche Lösung am besten zu eurer Geschichte und eurer Gruppe passt.

Es geht bei Heldenpfade nicht darum zu gewinnen oder zu verlieren. Das Ziel ist es, gemeinsam eine spannende Geschichte zu erleben.

In Heldenpfade erlebt ihr gemeinsam Abenteuer.

Ihr erkundet geheimnisvolle Ruinen, kämpft gegen gefährliche Gegner, helft Bewohnern eines Dorfes, entdeckt verborgene Schätze und erzählt gemeinsam eure eigene Geschichte.

Anders als bei vielen anderen Spielen gibt es dabei keinen Gewinner und keinen Verlierer. Statt gegeneinander zu spielen, arbeitet ihr als Gruppe zusammen und versucht gemeinsam, die Herausforderungen eures Abenteuers zu meistern.

Dieses Kapitel erklärt die Grundlagen von Heldenpfade und zeigt euch, wie ihr sofort mit eurem ersten Abenteuer beginnen könnt.

\section{Was ist Heldenpfade?}

Heldenpfade ist ein Pen-\&-Paper-Rollenspiel.

Gemeinsam erzählt ihr eine Geschichte und erlebt darin ein Abenteuer. Ein Spieler übernimmt die Spielleitung und erzählt das Abenteuer. Alle anderen Spieler übernehmen jeweils einen Helden und entscheiden selbst, was ihre Helden denken, sagen und tun.

Die Spieler sprechen miteinander, planen gemeinsam und versuchen als Gruppe, die Herausforderungen ihres Abenteuers zu meistern.

Die Spielleitung beschreibt die Welt, übernimmt alle anderen Figuren und entscheidet, wie die Umwelt auf die Handlungen der Helden reagiert. Sie spielt dabei nicht gegen die Spieler. Ihre Aufgabe ist es, gemeinsam mit allen am Tisch eine spannende Geschichte zu erzählen und dafür zu sorgen, dass alle Spaß am Spiel haben.

Dabei gibt es fast keine festen Grenzen. Vielleicht möchtet ihr mit einem Händler verhandeln, einen steilen Felsen erklimmen, ein Rätsel lösen oder gegen einen Drachen kämpfen. Fast jede Idee ist erlaubt.

Ob eine Handlung gelingt, entscheiden meistens die Regeln und die Würfel. Sie sorgen dafür, dass Erfolg und Misserfolg spannend bleiben und niemand vorher weiß, wie die Geschichte ausgeht.

Die Regeln sollen eure Fantasie dabei unterstützen -- sie sollen sie nicht einschränken.

\section{Wie läuft ein Spiel ab?}

Eine Partie Heldenpfade besteht aus einem ständigen Wechsel zwischen Erzählen, Entscheiden und Würfeln.

Zunächst beschreibt die Spielleitung eine Situation. Danach entscheiden die Spieler, wie ihre Helden darauf reagieren.

Ist unklar, ob eine Handlung gelingt, werden Würfel verwendet. Das Ergebnis entscheidet, wie die Geschichte weitergeht.

Ein kurzer Ausschnitt könnte zum Beispiel so aussehen:

\begin{quote}
\textbf{Spielleitung:} „Vor euch steht eine alte Holztür. Dahinter hört ihr leise Stimmen.“

\textbf{Spieler:} „Ich lausche zuerst an der Tür.“

\textbf{Spielleitung:} „Dann führe bitte eine Wahrnehmungsprobe durch.“
\end{quote}

Der Spieler würfelt und erfährt anschließend, was sein Held hört.

So entwickelt sich die Geschichte Schritt für Schritt weiter.

\section{Kann ich sofort losspielen?}

Ja.

Heldenpfade enthält mehrere spielfertige Beispielhelden. Ihr könnt euch jeweils einen dieser Helden aussuchen und sofort mit eurem ersten Abenteuer beginnen.

Die Beispielhelden befinden sich in Anhang 5 -- Beispielhelden.

Wenn ihr lieber einen eigenen Helden erschaffen möchtet, findet ihr die vollständigen Regeln dazu in Kapitel 9 -- Wie entsteht ein Held?

\section{Was brauche ich zum Spielen?}

Zum Spielen von Heldenpfade benötigt ihr:

\begin{itemize}
  \item dieses Regelwerk,
  \item ein Abenteuer,
  \item Charakterbögen mit Beispielhelden,
  \item Karten,
  \item Marker,
  \item Stifte
  \item sowie die zum Spiel gehörenden Würfel.
\end{itemize}

\section{Welche Würfel gibt es?}

In Heldenpfade kommen verschiedene Würfelarten zum Einsatz. Die Symbole werden im gesamten Regelwerk immer gleich verwendet.

\begin{center}
\begin{hptable}{C{1.7cm} L{4.0cm} Y}
\textbf{Symbol} & \textbf{Würfelart} & \textbf{Verwendung} \\
\heroDie{} & Heldenwürfel & Der persönliche Würfel eines Helden. Jeder Held besitzt genau einen Heldenwürfel. \\
\greenDie{} & Attributswürfel & Grüne Würfel für die natürlichen Eigenschaften und grundlegenden Stärken eines Helden. \\
\blueDie{} & Fertigkeitswürfel & Blaue Würfel für erlernte, geübte und trainierte Fähigkeiten. \\
\orangeDmgDie{} & Leichter Schadenswürfel & Orangener Würfel für leichten Schaden. \\
\redDmgDie{} & Mittlerer Schadenswürfel & Roter Würfel für mittleren Schaden. \\
\blackDmgDie{} & Schwerer Schadenswürfel & Schwarzer Würfel für schweren Schaden. \\
\end{hptable}
\end{center}

Welche Würfel bei einer Probe verwendet werden, hängt von den Fähigkeiten des Helden ab. Je gefährlicher ein Angriff ist, desto stärkere Schadenswürfel können zum Einsatz kommen.

Die genaue Verwendung aller Würfel wird in den folgenden Kapiteln erklärt.

\section{Wie ist dieses Regelwerk aufgebaut?}

Das Regelwerk ist so aufgebaut, dass ihr die Spielregeln Schritt für Schritt kennenlernt.

Die ersten Kapitel erklären die wichtigsten Grundlagen des Spiels. Danach folgen die Regeln für Kampf, Ausrüstung, Zauber und Talente. Anschließend erfahrt ihr, wie ihr euren eigenen Helden erschafft und weiterentwickelt.

Wenn ihr später eine bestimmte Regel nachschlagen möchtet, könnt ihr direkt zum passenden Kapitel springen.
